\documentclass{vgtc}                          % final (conference style)

\graphicspath{{figures/}{pictures/}{images/}{./}}

\usepackage{times}
\renewcommand*\ttdefault{txtt}
\usepackage{booktabs}
\usepackage{mathptmx}

\onlineid{0}
\vgtccategory{Research}
\vgtcinsertpkg

%% Paper title.
\title{PhaseViz: A Phase-Aware Post-Match Soccer Analytics Dashboard}

%% Author and Affiliation (single author).
\author{Omar A.\thanks{e-mail: placeholder@uwaterloo.ca}\\
        \scriptsize University of Waterloo}

%% Teaser figure
\teaser{
  \centering
  % TODO: Replace with actual screenshot
  % \includegraphics[width=\linewidth]{teaser}
  \caption{PhaseViz analysis view showing coordinated visualizations: phase detection timelines, pitch canvas with shape graph overlay and jersey-numbered players, collapsible team metric panels, and game momentum chart with goal markers. The scrubber bar features distinct event markers for goals, high-xG shots, and high-xT passes.}
  \label{fig:teaser}
}

%% Abstract
\abstract{
Post-match tactical analysis in professional soccer requires analysts to manually cross-reference video, event logs, and statistical reports to understand how teams transition between tactical phases such as pressing, build-up play, and counter-attacks. We present PhaseViz, an interactive dashboard that automatically detects tactical phases from 25\,Hz optical tracking data and renders coordinated visualizations linking phase context to spatial, temporal, and statistical views. A Python pipeline processes Sportec/DFL tracking and event data through an eight-step computation chain---feature extraction, two-path phase classification, formation detection, pitch control estimation, and expected threat modeling---producing static JSON files consumed by a React and D3.js frontend. The dashboard provides phase-annotated timelines, an animated pitch canvas with multiple overlay modes, diverging bar charts for game momentum and threat flow, and real-time team metrics. Brushing and linking across all views enables analysts to explore how spatial configurations, formation structure, and threat generation vary across different tactical phases.
}

%% Keywords
\keywords{Soccer analytics, tactical phases, coordinated views, tracking data, information visualization.}

%%%%%%%%%%%%%%%%%%%%%%%%%%%%%%%%%%%%%%%%%%%%%%%%%%%%%%%%%%%%%%%%
\begin{document}

\firstsection{Introduction}

\maketitle

Professional soccer analysis increasingly relies on spatiotemporal tracking data captured at 25\,Hz by optical systems. While raw tracking data provides precise positional information for all 22 players and the ball, extracting tactical meaning from millions of coordinate pairs remains challenging. Coaches and performance analysts typically rely on manual video review augmented by basic statistical summaries, a process that is time-consuming and difficult to scale across multiple matches.

A fundamental gap in current analytics tooling is the lack of \emph{phase awareness}. Soccer matches unfold as sequences of tactical phases---high pressing, defensive blocking, build-up play, counter-attacks---each with distinct spatial signatures and strategic implications. Existing dashboards either present raw tracking animations without phase context or show aggregate statistics that obscure within-match temporal dynamics.

We present \textbf{PhaseViz}, a post-match analytics dashboard that bridges this gap through three contributions:
\begin{enumerate}
  \item A \textbf{two-path phase classification algorithm} that combines tracking-based features (defensive line height, compactness, pressure proximity) with event-based counter-attack detection rules to segment matches into seven tactical phase types.
  \item A \textbf{coordinated multi-view visualization} system linking phase-annotated timelines, an animated pitch canvas with selectable overlays, diverging momentum charts, and live team metric panels through brushing and linking.
  \item A \textbf{static deployment architecture} where a Python pipeline preprocesses tracking data into JSON files consumed by a client-side React application, enabling deployment without a backend server.
\end{enumerate}

% ============================================================
\section{Background \& Related Work}
% ============================================================

\subsection{Soccer Visualization Systems}

Several systems have explored soccer data visualization. Perin et al.~\cite{perin2013} introduced \emph{SoccerStories}, enabling interactive exploration and communication of soccer match incidents through visual narratives. Janetzko et al.~\cite{janetzko2014} proposed feature-driven visual analytics of soccer data using automated event detection. Stein et al.~\cite{stein2018} developed multi-perspective analysis combining video and movement data for team sport analysis.

Commercial platforms such as The Athletic's match dashboard~\cite{athletic2023} popularized momentum charts and comparative team statistics. However, these systems generally present pre-computed summaries without exposing the underlying tactical phase structure or linking temporal patterns to spatial configurations.

\subsection{Tactical Phase Detection}

Our phase classification draws on established methods. Bekkers and Sahasrabudhe~\cite{bekkers2023} define counter-attack identification rules based on possession recovery location, forward distance, and ball velocity. Bauer and Anzer~\cite{bauer2021} formalize pressing detection through tracking-derived features. Our system combines these approaches into a unified two-path classifier that simultaneously labels both teams per frame.

\subsection{Formation and Spatial Analysis}

For formation detection, we adapt the RoleRep algorithm~\cite{bialkowski2014}, which uses an EM procedure with Hungarian assignment to discover stable player roles from positional data. Shape graph construction follows Brandes et al.~\cite{brandes2025}, applying iterative Delaunay edge removal based on angular stability to reveal meaningful spatial relationships between players. Expected Threat (xT) values use Karun Singh's Markov chain model~\cite{singh2019}, providing a spatial threat surface for evaluating ball progression.

% ============================================================
\section{Design}
% ============================================================

\subsection{Data Abstraction}

PhaseViz consumes the Sportec/DFL IDSSE dataset~\cite{bassek2025}, comprising seven Bundesliga matches with synchronized 25\,Hz optical tracking and event data. The pipeline transforms this into six JSON files per match:

\begin{itemize}
  \item \textbf{Frames} (${\sim}$2,700 frames at 4\,Hz): Player positions (\emph{x, y, speed}), ball position, jersey numbers.
  \item \textbf{Phases} (${\sim}$1,100 segments): Phase type, team, start/end time, duration, xThreat gained/conceded.
  \item \textbf{Formations} (${\sim}$90 per match): Mean role positions ($10{\times}2$), adjacency matrices ($10{\times}10$), stability scores, formation labels.
  \item \textbf{Pitch control}: Per-frame control grids ($24{\times}16$) clipped to convex hull.
  \item \textbf{Heatmaps}: KDE-based pressing intensity grids ($21{\times}14$) per high-press phase.
  \item \textbf{Metadata}: Teams, players, goals, match statistics, club badges.
\end{itemize}

\subsection{Task Abstraction}

Through informal discussions with soccer analytics practitioners, we identified five core analytical tasks:

\begin{itemize}
  \item \textbf{T1 -- Discover}: Identify which tactical phases dominated each team's match strategy.
  \item \textbf{T2 -- Compare}: Compare spatial metrics (defensive line height, compactness, pressure) across phase types.
  \item \textbf{T3 -- Locate}: Find specific high-threat moments---goals, dangerous shots, progressive passes.
  \item \textbf{T4 -- Browse}: Scrub through match time to observe how spatial configurations evolve.
  \item \textbf{T5 -- Correlate}: Link phase type to formation structure, pitch control patterns, and threat generation.
\end{itemize}

\subsection{Audience and Usage Scenario}

The primary audience is performance analysts at professional clubs who conduct post-match tactical reviews. A typical scenario: an analyst loads a match, identifies that the opponent's high-press phases coincided with the team's highest xThreat moments, scrubs to those time windows to examine the spatial configuration, and notes that a particular shape graph during build-up phases left central channels vulnerable.

\subsection{Visual Encoding Design}

The dashboard employs multiple coordinated views, each optimized for a specific data type.

\textbf{Phase Detection Timeline.} Two horizontal timelines (one per team) encode phase segments as colored rectangles. The \emph{x-position} and \emph{width} channels encode start time and duration; the \emph{color} channel (hue) encodes phase type using a seven-color muted earth-tone palette designed for cohesion: olive green (attacking), sage (build-up), terracotta (high press), taupe (mid block), steel blue-gray (defensive block), muted gold (counter-attack), and warm gray (open play). Goal markers appear as soccer ball icons at the scoring minute. A vertical current-time indicator supports temporal orientation.

\textbf{Pitch Canvas.} A Canvas-rendered pitch view uses \emph{point marks} for players (Canvas rather than SVG for performance at 4\,Hz with 22 players). The \emph{x,y position} channels map normalized coordinates to pitch space; \emph{color} encodes team identity using each club's badge-derived brand color; jersey \emph{numbers} (bold text inside circles) provide player identification. Three selectable overlays provide spatial analysis:
\begin{itemize}
  \item \emph{Shape Graph}: Iterative Delaunay edge removal by angular stability~\cite{brandes2025}, encoding formation structure.
  \item \emph{Pitch Control}: Colored dots on a $24{\times}16$ grid clipped to the attacking team's convex hull.
  \item \emph{Clean}: Unadorned player positions.
\end{itemize}

High-xT passes appear as dashed team-colored arrows with pill-shaped value labels. High-xG shots appear as rings with radius proportional to xG, using filled team-colored backgrounds for goals.

\textbf{Game Momentum Chart.} A diverging bar chart with \emph{minute} on the x-axis. Home team momentum extends upward; away team downward. Momentum is computed from per-event xThreat blended with shot xG, smoothed via bidirectional exponential moving average. Goal markers (soccer ball icons) appear above or below the chart at the scoring minute.

\textbf{Team Metric Panels.} Two collapsible side panels (one per team, horizontally expanding) display real-time metrics: current phase label, defensive line height, compactness, pressure proxy, team length, and team width. Phase distribution appears as colored dot strips proportional to time in each phase type.

\textbf{Playback Scrubber.} A range slider with event markers overlaid using distinct shapes: \emph{circles} for goals, \emph{rounded squares} for high-xG shots, and \emph{triangles} for high-xT passes, all in team colors. A legend identifies each shape.

\subsection{Interaction Design}

Following Shneiderman's visual information-seeking mantra~\cite{shneiderman1996}:

\textbf{Overview first}: The full match is visible across the phase timeline, momentum chart, and playback scrubber simultaneously.

\textbf{Zoom and filter}: Phase type toggles filter timelines to show only selected phases. A team selector switches between home, away, or both teams.

\textbf{Details on demand}: Hovering over phase segments reveals tooltips with type and duration. The scrubber supports click-to-jump and drag-to-scrub. Playback speed is adjustable ($0.5{\times}$ to $4{\times}$).

\textbf{Brushing and linking}: Scrubbing updates the pitch canvas, chart time indicators, side panel metrics, and live score simultaneously. Switching overlay modes preserves the current time position across all views.

% ============================================================
\section{Implementation}
% ============================================================

The system comprises a Python preprocessing pipeline and a React frontend.

\textbf{Pipeline.} Nine scripts process each match sequentially: (1)~load Sportec XML via kloppy~\cite{kloppy}; (2)~compute sliding-window features (defensive line height from the mean x of the four deepest outfield defenders, convex hull area for compactness, and a count of defenders within 5m of the ball carrier for pressure proxy); (3)~classify phases via a two-path algorithm---Path A applies tracking-based thresholds (e.g., defensive line height $>$ 0.38 and ball in opponent half for high press, compactness $<$ 0.20 and deep defensive line for block), while Path B detects counter-attacks using event-based rules~\cite{bekkers2023} (recovery in own half, $>$7.5m forward, $>$3 m/s velocity); phases shorter than 5 seconds are discarded; (4)~estimate pitch control on a $24{\times}16$ grid within the attacking team's convex hull; (5)~detect formations using EM iteration with Hungarian assignment~\cite{bialkowski2014} and shape graph edge construction~\cite{brandes2025}; (6)~compute pressing heatmaps via 2D KDE at $21{\times}14$ resolution; (7)~compute per-event xThreat using a $12{\times}8$ Markov chain grid~\cite{singh2019}, with special handling for goals (xT = 0.50) and shots (positional xG model); (8)~downsample tracking from 25\,Hz to 4\,Hz and export six JSON files; (9)~compute match-level and player-level statistics (possession, shots, progressive passes, expected goals).

\textbf{Frontend.} Built with React~18 and D3.js~7. The pitch canvas uses the HTML5 Canvas API for smooth rendering of 22 animated player dots---SVG would create ${\sim}$22 DOM nodes per frame, degrading performance during playback. All charts use D3 rendering into SVG. The application loads static JSON files from the public directory, requiring no backend server. Team colors and badges are sourced from a centralized match catalog for instant match switching. The design adopts The Athletic's match dashboard aesthetic~\cite{athletic2023}: a flat layout on warm cream background with muted team-derived colors, compact typography (Inter font family), and thin divider lines between sections rather than card containers.

% ============================================================
\section{Proposed Evaluation}
% ============================================================

We propose three evaluation strategies:

\textbf{Expert review.} Present the dashboard to 2--3 professional soccer analysts or coaching staff. Using a think-aloud protocol, ask them to: (a)~identify the three most dangerous counter-attacks, (b)~compare pressing intensity between halves, and (c)~assess whether formation changes correlate with phase transitions. Measure task completion time and qualitative feedback.

\textbf{Comparative task study.} Compare PhaseViz against a baseline of video review with basic event logs. Give participants (graduate students with soccer knowledge) a set of tactical questions and measure time-to-insight and answer accuracy.

\textbf{Phase classification validation.} Sample 50 phase segments, manually annotate them from match video, and compute agreement with the automated classifier. This validates the pipeline's accuracy, which directly affects the dashboard's utility.

% ============================================================
\section{Discussion \& Future Work}
% ============================================================

\subsection{Lessons Learned}

\textbf{Phase thresholds require domain calibration.} Our tracking-based thresholds were calibrated iteratively by inspecting the dashboard output against match video. A supervised learning approach trained on analyst annotations would be more robust but would require labeled training data that is not publicly available.

\textbf{Canvas outperforms SVG for animated spatial views.} Rendering 22 player circles with jersey numbers at 4\,Hz was infeasible with SVG due to DOM manipulation overhead. Canvas provided smooth playback even with shape graph edges and pitch control overlays rendered simultaneously.

\textbf{Static deployment enables broad access.} By preprocessing all data into JSON files, the dashboard runs as a static site without backend infrastructure. This simplifies deployment and removes latency from data queries, at the cost of requiring the full pipeline to be rerun when parameters change.

\subsection{Limitations}

The phase classifier uses empirically tuned thresholds rather than learned parameters, limiting generalizability across leagues or playing styles. The xThreat model uses a pre-trained $12{\times}8$ grid rather than a model trained on the specific league's data. Only one match has been fully demonstrated in the current prototype, though the pipeline supports all seven IDSSE matches. No formal user evaluation has been conducted.

\subsection{Future Work}

Future directions include: (1)~training a machine learning phase classifier on analyst-annotated data; (2)~cross-match comparison views for the same team across different opponents; (3)~player-level phase contribution analysis; (4)~physics-based pitch control models following Spearman~\cite{spearman2018}; and (5)~real-time streaming integration for live match analysis.

% ============================================================
\section{Conclusion}
% ============================================================

PhaseViz demonstrates that automated tactical phase detection, combined with coordinated multi-view visualization, can provide analysts with a structured lens for post-match review. By segmenting matches into semantically meaningful phases and linking them to spatial, temporal, and statistical views, the dashboard enables exploration of tactical patterns that would otherwise require extensive manual video review. The static deployment architecture and support for multiple matches make it a practical foundation for soccer analytics workflows.

% ============================================================
% References (up to 1 extra page allowed)
% ============================================================

\bibliographystyle{abbrv-doi}

\begin{thebibliography}{99}

\bibitem{bassek2025}
M.~Bassek, et al.
\newblock An integrated dataset of spatiotemporal and event data in elite soccer.
\newblock {\em Scientific Data}, 12(1), 2025.

\bibitem{bauer2021}
P.~Bauer and G.~Anzer.
\newblock Data-driven detection of counterpressing in professional football.
\newblock {\em Data Mining and Knowledge Discovery}, 35(5):2009--2049, 2021.

\bibitem{bekkers2023}
J.~Bekkers and A.~Sahasrabudhe.
\newblock Identifying counterattacks in football.
\newblock In {\em MIT Sloan Sports Analytics Conference}, 2023.

\bibitem{bialkowski2014}
A.~Bialkowski, P.~Lucey, P.~Carr, Y.~Yue, S.~Sridharan, and I.~Matthews.
\newblock Identifying team style in soccer using formations learned from spatiotemporal tracking data.
\newblock In {\em IEEE ICDM Workshops}, pp.~9--14, 2014.

\bibitem{brandes2025}
U.~Brandes, J.~Sotudeh, and P.~Borm.
\newblock Shape graphs for soccer formation analysis.
\newblock Technical report, ETH Zurich, 2025.

\bibitem{janetzko2014}
H.~Janetzko, D.~Sacha, M.~Stein, T.~Schreck, D.~Keim, and O.~Deussen.
\newblock Feature-driven visual analytics of soccer data.
\newblock In {\em IEEE VAST}, pp.~13--22, 2014.

\bibitem{kloppy}
kloppy contributors.
\newblock kloppy: Standardizing soccer tracking and event data.
\newblock \url{https://github.com/PySport/kloppy}, 2024.

\bibitem{munzner2014}
T.~Munzner.
\newblock {\em Visualization Analysis and Design}.
\newblock CRC Press, 2014.

\bibitem{perin2013}
C.~Perin, R.~Vuillemot, and J.-D.~Fekete.
\newblock {S}occer{S}tories: A kick-off for visual soccer analysis.
\newblock {\em IEEE Trans.\ Vis.\ Comput.\ Graphics}, 19(12):2506--2515, 2013.

\bibitem{shneiderman1996}
B.~Shneiderman.
\newblock The eyes have it: A task by data type taxonomy for information visualizations.
\newblock In {\em IEEE Symp.\ Visual Languages}, pp.~336--343, 1996.

\bibitem{singh2019}
K.~Singh.
\newblock Introducing expected threat ({xT}).
\newblock \url{https://karun.in/blog/expected-threat.html}, 2019.

\bibitem{spearman2018}
W.~Spearman.
\newblock Beyond expected goals.
\newblock In {\em MIT Sloan Sports Analytics Conference}, 2018.

\bibitem{stein2018}
M.~Stein, H.~Janetzko, et al.
\newblock Bring it to the pitch: Combining video and movement data to enhance team sport analysis.
\newblock {\em IEEE Trans.\ Vis.\ Comput.\ Graphics}, 24(1):13--22, 2018.

\bibitem{athletic2023}
The Athletic.
\newblock Match dashboard.
\newblock \url{https://www.nytimes.com/athletic/}, 2023.

\end{thebibliography}

\end{document}
